\section{Background}
\label{sec:background}

Both deployed FP4 formats store identical E2M1 elements (largest finite magnitude
6) and differ only in how a block is scaled (Table~\ref{tab:formats}).

\begin{table}[t]
  \centering
  \caption{The two block-scaled FP4 formats. Elements are identical; the scaling
    is the whole of the difference.}
  \label{tab:formats}
  \begin{tabular}{@{}l p{3.4cm} p{4.2cm}@{}}
    \toprule
                              & MXFP4                          & NVFP4 \\
    \midrule
    element format            & E2M1 (1/2/1)                    & E2M1 (1/2/1) \\
    block size                & 32                             & 16 \\
    block scale format        & E8M0 (0/8/0), a power of two    & E4M3 (1/4/3) \\
    global scale              & none                           & one per tensor (FP32) \\
    scale range               & $2^{-127}$ to $2^{127}$ (255 binades) & $2^{-9}$ to 448 relative to the global scale ($\sim$19 binades) \\
    scale resolution          & one binade (no mantissa bits)  & 6.25\% (3 mantissa bits) \\
    clipping of the block max  & impossible by construction     & up to $\sim$6\% \\
    specification / vendor    & OCP MX v1.0 \citep{ocp2023mx}   & NVIDIA \citep{nvidia2026nvfp4} \\
    \bottomrule
  \end{tabular}

  \vspace{0.6em}
  {\footnotesize Blocks run along the contraction axis. The MXFP4 block-scale
  exponent admits more than one rounding convention in the literature; the one
  used here is stated in Section~\ref{sec:method}. Scale-resolution and clipping
  figures follow from the format definitions, not derived here.}
\end{table}

The classical tools for this kind of product are the rounding-error bounds for an
inner product. For length $n$ and unit roundoff $u$, the deterministic worst case
is
$\gamma_n = n\,u/(1-n\,u)$ \citep{higham2002accuracy}, obtained by compounding one
rounding per accumulation step. A randomized analysis gives a more refined
estimate for typical cases, changing the leading factor to give a high-probability
bound of order $\sqrt{n \log n}\,u$ \citep{higham2019new}. Both are stated in
terms of a single unit roundoff $u$ fixed by the encoding.

NVIDIA's reported 36\%-more-tokens gap between the two formats
\citep{nvidia2026nvfp4} is given without an explanation of where the difference
comes from.
